% !TeX root = ../../Book4.tex
% 纲要标题：### 2.4 显式覆盖的几何
% 纲要位置：几何学计划纲要.md，第 3536 行。
% 正文状态：初稿；尚未正式审核。

% 纲要内容开始
% 
% * 超椭圆曲面的分支粘合
% * 三次曲线及亏格
% * 拓扑单值与代数方程的对应
% * 不可约、连通与分支数据
% 
% ---
% 
% 纲要内容结束

\section{显式覆盖的几何}
\label{g4:sec:0204}

方程 \(y^2=P(z)\) 在一个普通的 \(z\) 值上有两个根。沿闭路继续这两个根，返回时可能交换；在 \(P\) 的零点处，两根则合并为一个分歧点。把这些现象组织成一张曲面，需要同时回答三个问题：局部怎样给合并点取坐标，无穷远处应当补几个点，以及所有局部分支究竟连成几张连通曲面。下面先在二次方程中逐项完成构造，再说明一般代数方程中的根置换如何回答连通性问题。

\subsection{超椭圆曲面的分支粘合}
\label{g4:subsec:0204:01}

设
\[
	P(z)=a_d\prod_{j=1}^d(z-a_j),\qquad a_d\ne0,\quad d\ge1,
\]
其中 \(a_1,\ldots,a_d\) 两两不同。先取仿射方程的解集
\[
	X_P^\circ=\{(z,y)\in\mathbb C^2:y^2=P(z)\},
	\qquad \pi(z,y)=z.
\]
在 \(P(z)\ne0\) 处，可以在小圆盘内选取平方根，解集分成两个全纯图像。根相遇的地方要换用 \(y\) 作局部参数；无穷远则必须根据 \(P\) 的次数另取坐标。

\begin{proposition}\label{g4:prop:0204:hyperelliptic-completion}
	设 \(P\in\mathbb C[z]\) 为次数 \(d\ge1\) 的无重根多项式。上述 \(X_P^\circ\) 是连通黎曼曲面，并可补成紧黎曼曲面 \(X_P\)，使 \(\pi\) 延拓为二次全纯映射
	\[
		\pi:X_P\longrightarrow\mathbb P^1.
	\]
	每个有限零点 \(a_j\) 上方有一个分歧指数为 \(2\) 的点。当 \(d\) 为偶数时，无穷远上方有两个无分歧点；当 \(d\) 为奇数时，无穷远上方只有一个点，分歧指数为 \(2\)。其余各点均无分歧。
\end{proposition}
\begin{proof}
	先构造有限处的坐标。若 \(P(z_0)\ne0\)，在 \(z_0\) 的小圆盘上选取全纯平方根 \(s(z)\)，两张图像 \(y=\pm s(z)\) 都以 \(z\) 为坐标，投影没有分歧。若 \(P(a)=0\)，则 \(P'(a)\ne0\)，由一元全纯逆函数定理，\(P\) 在 \(a\) 附近有全纯逆函数 \(\psi\)。因此原方程在 \((a,0)\) 附近写成
	\[
		z=\psi(y^2),
	\]
	并以 \(y\) 为坐标。若进一步写 \(P(z)=(z-a)U(z)\)，在小圆盘上取 \(q(z)^2=U(z)\)，则
	\[
		t=\frac{y}{q(z)},\qquad z-a=t^2
	\]
	是把投影化成二次幂的坐标。这里 \(q(a)\ne0\)，所以 \(t\) 对 \(y\) 的导数不为零。这些图册的过渡由方程中的全纯函数给出。

	令
	\[
		Q(u)=u^dP(1/u)=a_d+a_{d-1}u+\cdots+a_0u^d,
	\]
	于是 \(Q(0)=a_d\ne0\)。取原点附近的全纯平方根 \(q(u)^2=Q(u)\)。

	当 \(d=2m\) 时，在 \(|z|\) 很大处使用
	\[
		u=\frac1z,\qquad v=y\,u^m.
	\]
	方程化为 \(v^2=Q(u)\)。这给出两张图像 \(v=q(u)\) 与 \(v=-q(u)\)，各补一个点
	\[
		p_\infty^+=(0,q(0)),\qquad p_\infty^-=(0,-q(0)).
	\]
	两处都以 \(u\) 为坐标；目标球面在无穷远的坐标同样是 \(u=1/z\)，故投影在两处均无分歧。

	当 \(d=2m+1\) 时，改用
	\[
		u=\frac1z,\qquad v=y\,u^{m+1}.
	\]
	方程化为 \(v^2=uQ(u)\)，在 \(u=0\) 上只有 \((0,0)\)。令 \(t=v/q(u)\)，便得到
	\[
		u=t^2,\qquad v=tq(t^2).
	\]
	这是一张以 \(t\) 为坐标的圆盘，去掉 \(t=0\) 后与原曲面在无穷远附近的部分对应。只补点 \(p_\infty\) 即可，且投影在目标坐标 \(u\) 中为 \(u=t^2\)。在这张坐标中，原函数为
	\begin{equation}\label{g4:eq:0204:odd-infinity}
		z=t^{-2},\qquad y=t^{-(2m+1)}q(t^2).
	\end{equation}

	按照这些坐标与原曲面的公共部分粘合，得到无边界的一维复流形。它是 Hausdorff 的：投影不同的点可由球面中的不交邻域分开；同一有限纤维中的不同点由两张平方根图像分开；偶数次数时新补的两个点有不交的坐标圆盘。原仿射部分第二可数，新增图册只有一张或两张，故整个空间第二可数。

	紧性可以直接检查。对充分大的 \(R\)，集合
	\[
		K_R=\{(z,y)\in X_P^\circ:|z|\le R\}
	\]
	是 \(\mathbb C^2\) 中的闭有界集，因为 \(|y|^2=|P(z)|\) 在 \(|z|\le R\) 上有统一上界，所以 \(K_R\) 紧。无穷远附近的一张或两张闭坐标圆盘覆盖 \(|z|\ge R\) 的部分，并包含所有新补点。因此 \(X_P\) 是有限个紧集的并。

	最后说明连通性。去掉有限分支值后，底空间 \(\mathbb C\setminus\{a_1,\ldots,a_d\}\) 道路连通，投影是二叶覆盖。绕任一 \(a_j\) 的小正向圆周一周，局部公式 \(z-a_j=t^2\) 使两张覆盖片交换。因此基点纤维中的两个点可沿提升道路相连；其他纤维中的点也可沿底空间道路提升连到它们。这个无分歧部分连通，所有分歧点和新补点均在其闭包中，故 \(X_P^\circ\) 和 \(X_P\) 都连通。
\end{proof}

这个构造也可用剪开后跨缝粘合来表示。旋转 \(z\) 平面，使各 \(a_j\) 的实部互不相同，并按实部递增排列。依次用线段连接 \(a_1,a_2\)，连接 \(a_3,a_4\)，如此继续；这些线段的实部区间不交，因而彼此不交。如果 \(d\) 为奇数，再从最后的 \(a_d\) 向实部增加的方向画一条到无穷远的水平射线。删去这些切缝所得的开集 \(U\) 道路连通：连接两点的折线若遇到切缝，可以沿切缝的一侧绕过有限端点，作有限次小的绕行后便避开全部切缝。

\ProseParagraphBreak

在 \(U\) 上可以选取单值全纯平方根。为核对这一点，取 \(U\) 中任一闭路 \(\gamma\)。同一有限切缝的两个端点沿切缝相连而不遇到 \(\gamma\)，所以它们关于 \(\gamma\) 的绕数相等；若有通向无穷远的射线，沿射线绕数恒定，并在足够远处为零，故该端点的绕数为零。于是
\[
	\frac1{2\pi\mathrm{i}}\int_\gamma\frac{P'(z)}{P(z)}\,\mathrm{d}z
	=\sum_{j=1}^d\operatorname{Ind}(\gamma,a_j)
\]
是偶数。固定 \(z_0\in U\) 及 \(c^2=P(z_0)\)，公式
\[
	s(z)=c\exp\left(\frac12\int_{z_0}^z\frac{P'(\zeta)}{P(\zeta)}
	\,\mathrm{d}\zeta\right)
\]
因此与 \(U\) 中的路径选择无关，并满足 \(s^2=P\)。

\ProseParagraphBreak

取两张 \(U\)，分别标记 \(y=s(z)\) 与 \(y=-s(z)\)。跨过一条切缝时，平方根的两种边界值互为相反数：这在端点附近由 \(t^2=z-a_j\) 看出，并沿切缝内部由解析延拓保持。因此粘合规则为
\[
	\begin{array}{c@{\ \longleftrightarrow\ }c}
		\text{第一张的上岸}&\text{第二张的下岸}\\
		\text{第一张的下岸}&\text{第二张的上岸}.
	\end{array}
\]
这里上下岸指沿切缝选定的两侧，粘合保持同一个 \(z\) 值和同一个 \(y\) 值。缝的端点用 \(t\) 圆盘补入，便没有留下边界；无穷远端的补法由次数奇偶性决定。这种粘合描述与方程构造给出同一个带投影的曲面。

\begin{corollary}\label{g4:cor:0204:hyperelliptic-genus}
	设 \(P\in\mathbb C[z]\) 无重根，且 \(\deg P=d\ge1\)。上述紧曲面 \(X_P\) 的亏格为
	\[
		g(X_P)=\left\lfloor\frac{d-1}{2}\right\rfloor.
	\]
\end{corollary}
\begin{proof}
	映射 \(\pi:X_P\to\mathbb P^1\) 的次数为 \(2\)，目标亏格为 \(0\)。每个分歧点的指数为 \(2\)，所以它对分歧除子的次数贡献 \(1\)。由\cref{g4:thm:0203:riemann-hurwitz}，
	\[
		2g(X_P)-2=
		\begin{cases}
			-4+d,&d\text{ 为偶数},\\
			-4+(d+1),&d\text{ 为奇数}.
		\end{cases}
	\]
	分别解出 \(g=(d-2)/2\) 与 \(g=(d-1)/2\) 即得结论。
\end{proof}

亏格至少为 \(2\)，并存在到 \(\mathbb P^1\) 的二次全纯映射的紧黎曼曲面，称为\term[chaotuoyuanqumian]{超椭圆曲面}[hyperelliptic Riemann surface]。上述方程在 \(d\ge5\) 时给出这样的曲面；\(d=1,2\) 时得到亏格 \(0\) 的二叶模型，\(d=3,4\) 时得到亏格 \(1\) 的模型。无重根条件使有限处的合并都恰为 \(t\mapsto t^2\)，并保证仿射方程本身已经光滑。

\subsection{三次曲线及亏格}
\label{g4:subsec:0204:02}

三次方程把上面的奇数次数模型放进普通射影平面，不需要另选加权坐标。取
\[
	P(z)=a_3z^3+a_2z^2+a_1z+a_0,\qquad a_3\ne0,
\]
并假设 \(P\) 无重根。在齐次坐标 \([Z:Y:T]\) 中，考虑射影三次曲线
\begin{equation}\label{g4:eq:0204:projective-cubic}
	Y^2T=a_3Z^3+a_2Z^2T+a_1ZT^2+a_0T^3.
\end{equation}
在 \(T\ne0\) 的仿射图中，\(z=Z/T\)、\(y=Y/T\)，方程恢复为 \(y^2=P(z)\)。

\begin{proposition}\label{g4:prop:0204:cubic}
	设三次多项式 \(P\in\mathbb C[z]\) 无重根。方程\eqref{g4:eq:0204:projective-cubic}定义一条光滑射影三次曲线 \(E\)，它与\cref{g4:prop:0204:hyperelliptic-completion}构造的 \(X_P\) 双全纯。无穷远直线 \(T=0\) 与它只有一个交点 \(O=[0:1:0]\)，且
	\[
		g(E)=1.
	\]
	仿射部分的微分 \(\omega=\mathrm{d}z/y\) 延拓为 \(E\) 上处处非零的全纯微分。
\end{proposition}
\begin{proof}
	在仿射图中，方程函数 \(F(z,y)=y^2-P(z)\) 的两个偏导数为 \(-P'(z)\) 与 \(2y\)。它们不能同时为零，否则 \(P\) 有重根；有限处的全纯坐标已经在前一命题中给出。令 \(T=0\)，齐次方程迫使 \(Z=0\)，于是无穷远只有 \(O=[0:1:0]\)。在齐次方程左减右所得多项式中，对 \(T\) 的偏导数在 \(O\) 为 \(1\)，所以这里也不出现奇点。

	还可直接写出 \(O\) 附近的坐标，核对它就是前面补入的点。令
	\[
		Q(u)=a_3+a_2u+a_1u^2+a_0u^3,\qquad q(u)^2=Q(u),\quad q(0)\ne0.
	\]
	由\eqref{g4:eq:0204:odd-infinity}，仿射部分在无穷远附近表示为
	\[
		z=t^{-2},\qquad y=t^{-3}q(t^2).
	\]
	相应射影点为
	\[
		[Z:Y:T]=[t:q(t^2):t^3].
	\]
	它在 \(t=0\) 取值 \(O\)。在 \(Y\ne0\) 的射影图中，第一坐标 \(Z/Y=t/q(t^2)\) 对 \(t\) 的导数在 \(0\) 为 \(1/q(0)\ne0\)，所以 \(t\) 确为这条射影曲线上的全纯局部坐标。仿射恒等映射因而在补点处也成为双全纯对应。曲面紧性由前一构造得到，亏格由\cref{g4:cor:0204:hyperelliptic-genus}取 \(d=3\) 得到。等价地，投影 \(z:E\to\mathbb P^1\) 在三个有限根和 \(O\) 处各有一次简单分歧，Hurwitz 公式给出 \(2g-2=-4+4=0\)。

	现在逐点检查 \(\omega\)。在 \(y\ne0\) 的有限点处，\(z\) 是坐标，系数 \(1/y\) 全纯且非零。在有限分歧点 \((a,0)\) 附近，\(y\) 是坐标；对方程求微分，得到
	\[
		2y\,\mathrm{d}y=P'(z)\,\mathrm{d}z,
		\qquad
		\omega=\frac{2\,\mathrm{d}y}{P'(z)}.
	\]
	由于 \(P'(a)\ne0\)，这个表达式延拓且不为零。最后，在 \(O\) 的 \(t\) 坐标中，
	\[
		\omega
		=\frac{-2t^{-3}\,\mathrm{d}t}{t^{-3}q(t^2)}
		=-\frac{2}{q(t^2)}\,\mathrm{d}t,
	\]
	仍全纯且非零。
\end{proof}

在一张复坐标图中，\(\mathrm{d}z/y\) 看起来在 \(y=0\) 处有分母为零的问题；换到真正的曲面坐标后，分子同时消去这一阶数。无穷远处也由同一次数计算消去了表面上的极点。因此 \(\operatorname{div}(\omega)=0\)，而 \(\omega\) 给出 \(E\) 的典范线丛的全局非消失截面。

\ProseParagraphBreak

紧黎曼曲面 \(E\) 的亏格为 \(1\)，并选定一点 \(O\in E\) 时，二者组成一条\term[tuoyuanquxian]{椭圆曲线}[elliptic curve]。上述三次模型以唯一的无穷远点作为 \(O\)。这里的方程、分歧点和非消失微分都已由局部计算确定；它们为随后研究积分与周期提供了具体对象。

\subsection{拓扑单值与代数方程的对应}
\label{g4:subsec:0204:03}

局部根的变化可以直接写成纤维上的置换。例如方程
\[
	w^n=z,\qquad n\ge1,
\]
在 \(z\in\mathbb C^\times\) 上定义 \(n\) 叶无分歧覆盖。固定正实基点 \(z_0\)，把它的根标为
\[
	w_j=z_0^{1/n}e^{2\pi\mathrm{i} j/n},\qquad j\in\mathbb Z/n\mathbb Z.
\]
沿 \(z=z_0e^{2\pi\mathrm{i} s}\)、\(0\le s\le1\) 继续时，根 \(w_j\) 变成 \(w_{j+1}\)，所以单值是一个 \(n\)-循环。用 \(w\) 作源坐标，整个紧化映射就是 \(\mathbb P^1\to\mathbb P^1\)、\(w\mapsto w^n\)。在无穷远处令 \(v=1/w\)、\(u=1/z\)，仍有 \(u=v^n\)；\(0\) 与无穷远的分歧指数都为 \(n\)。目标无穷远附近的正向小环路在 \(z\) 平面中沿相反方向绕行，故其单值是前一个循环的逆。

\ProseParagraphBreak

对 \(y^2=P(z)\)，绕每个简单有限零点的单值都是换位 \((12)\)。全部有限零点的环路依次作用，得到 \((12)^d\)，所以绕无穷远的单值在 \(d\) 偶数时为恒等，在 \(d\) 奇数时为换位。这恰好恢复了前面的两个补点与一个二重分歧补点。局部幂坐标、根的解析延拓和拓扑单值因此给出同一组数据。

\ProseParagraphBreak

现在固定球面上的有限非空集合 \(B=\{b_1,\ldots,b_r\}\) 及基点 \(z_0\notin B\)。在 \(U=\mathbb P^1\setminus B\) 中，选取从 \(z_0\) 出发、绕各 \(b_j\) 正向一周再返回的环路 \(\gamma_j\)，使
\[
	\pi_1(U,z_0)=
	\langle\gamma_1,\ldots,\gamma_r:
	\gamma_1*\cdots*\gamma_r=1\rangle.
\]
这个展示可从第一卷命题 5.2.10 直接得到：把 \(b_r\) 送到无穷远后，\(U\) 成为去掉 \(r-1\) 个点的平面，前 \(r-1\) 个绕点环路自由生成基本群；按切开路径的次序，它们的乘积是包围全部有限缺失点的大圆周，而无穷远的正向环路是其反向。这里仍约定 \(\alpha*\beta\) 先走 \(\alpha\) 后走 \(\beta\)。

\begin{theorem}\label{g4:thm:0204:monodromy-construction}
	固定 \(B,z_0,\gamma_1,\ldots,\gamma_r\) 如上，并取整数 \(n\ge1\)。连通紧黎曼曲面的 \(n\) 次全纯映射
	\[
		\pi:X\longrightarrow\mathbb P^1,
		\qquad \text{所有分支值均包含于 }B,
	\]
	在保持目标投影的双全纯等价下，与满足以下条件的置换元组 \((\sigma_1,\ldots,\sigma_r)\in S_n^r\) 在同时共轭下的等价类一一对应：
	\begin{enumerate}
		\item \(\sigma_r\circ\cdots\circ\sigma_1=\operatorname{id}\)；
		\item \(\sigma_1,\ldots,\sigma_r\) 生成的群在 \(\{1,\ldots,n\}\) 上传递。
	\end{enumerate}
	其中 \(\sigma_j\) 是沿 \(\gamma_j\) 提升后作用于基点纤维的终点置换。它的每条长度为 \(e\) 的循环对应 \(b_j\) 上方的一个点，该点的分歧指数为 \(e\)。这里同时共轭指用同一个 \(\tau\in S_n\) 把每个 \(\sigma_j\) 换成 \(\tau\sigma_j\tau^{-1}\)。
\end{theorem}
\begin{proof}
	先从给定映射得到数据。删去 \(B\) 及其原像，得到一个 \(n\) 叶无分歧覆盖。\cref{g4:thm:0202:unbranched-cover}的证明已说明，连通二维流形删去任意有限集后仍道路连通；这里 \(B\) 及其原像都有限，所以这个结论适用。因此基点纤维上的单值传递。环路先走 \(\gamma_1\) 再走 \(\gamma_2\) 时，终点置换为 \(\sigma_2\circ\sigma_1\)，故基本群的唯一关系给出第一项。局部幂模型说明，一个 \(e\) 次分歧点恰好使附近的 \(e\) 张片循环置换。重新编号纤维，正好产生同时共轭。

	反过来，满足第一项的元组给出 \(G=\pi_1(U,z_0)\) 在 \(S=\{1,\ldots,n\}\) 上的右作用，规则为 \(j\cdot\gamma_\nu=\sigma_\nu(j)\)。第一卷第5.5.3小节的覆盖重建定理 5.5.6 适用于 \(U\)：它道路连通，局部同胚于圆盘，所以局部道路连通且半局部单连通。具体地，取从 \(z_0\) 出发的道路类万有覆盖 \(\widetilde U\)，用
	\[
		g\cdot(j,[\alpha])=
		\bigl(j\cdot g^{-1},[\gamma*\alpha]\bigr),
		\qquad g=[\gamma],
	\]
	在 \(S\times\widetilde U\) 上取商，便得到具有指定单值的 \(n\) 叶覆盖 \(E\to U\)。该定理还说明，覆盖的等价恰好对应纤维的等变双射。每个连通分支都遇到基点纤维，而两个纤维点属于同一分支，当且仅当有底空间环路的提升连接它们；第二项因而保证 \(E\) 连通。

	有限叶性使 \(E\) 的覆盖片有可数基；同纤维的不同点由不同覆盖片分离，异纤维点由底空间分离。因此 \(E\) 是 Hausdorff 且第二可数的，\cref{g4:thm:0201:lift-complex-structure}给出唯一使投影局部双全纯的复结构。围绕每个 \(b_j\) 取互不相交的小坐标圆盘 \(D_j\)，把 \(b_j\) 的坐标记为 \(0\)。在穿孔圆盘 \(D_j^*\) 上，\cref{g4:lem:0202:punctured-disk-cover}把覆盖分成与 \(\sigma_j\) 的循环对应的有限个分支；长度为 \(e\) 的分支有坐标 \(t\)，投影写成 \(u=t^e\)。给每个这样的分支补入 \(t=0\)，便得到一张无边界复曲面 \(X\) 和全纯投影 \(\pi\)。原来的 \(E\) 稠密且连通，故 \(X\) 连通。

	所得曲面紧。确实，有限叶覆盖在底空间的每个紧集上有紧原像：用有限个闭小圆盘覆盖这个紧集，且各闭圆盘都包含在均匀覆盖的稍大圆盘内，它们的原像分别是有限个紧副本。球面去掉各 \(b_j\) 的更小圆盘后是紧集，其原像紧；再加上刚补入的有限个闭 \(t\) 圆盘，就覆盖整个 \(X\)。补点引理同时保证各新点有互不混淆的圆盘邻域，因而保留 Hausdorff 性和第二可数性。

	最后，两个单值元组只差同时共轭时，覆盖重建定理给出 \(E\) 之间的等价。它在每张覆盖片上由目标坐标表示为恒等，故全纯；由穿孔圆盘补点引理，它在每个新点处唯一延拓为双全纯映射。反过来，保持 \(\pi\) 的双全纯等价限制到基点纤维后与全部单值置换相容，必给出同时共轭。这也证明补点后没有额外的等价选择。
\end{proof}

定理允许 \(B\) 中有实际上不分歧的点：它们对应 \(\sigma_j=\operatorname{id}\)。如果要求分支值恰为 \(B\)，就要求每个 \(\sigma_j\) 非恒等。给定循环分解之后，补点的个数、每一点的指数以及对 Hurwitz 公式的贡献，都可以逐个读出。

\subsection{不可约、连通与分支数据}
\label{g4:subsec:0204:04}

对一般方程，连通性要在局部根组成的覆盖上判断。考虑关于 \(w\) 首一的多项式
\[
	P(z,w)=w^n+a_1(z)w^{n-1}+\cdots+a_n(z),
	\qquad a_j\in\mathbb C(z),\quad n\ge1.
\]
假设它在 \(\mathbb C(z)[w]\) 中可分，即 \(P\) 与 \(\partial P/\partial w\) 互素。互素性给出 \(A,B\in\mathbb C(z)[w]\) 使 \(AP+BP_w=1\)。清去这个恒等式中有理系数的分母后，说明只有有限多个 \(z\) 可能使专门化方程出现重根。因此可以选有限集合 \(S\subset\mathbb P^1\)，包含无穷远、全部系数极点及所有这样的例外值，使在 \(U=\mathbb P^1\setminus S\) 上，方程始终有 \(n\) 个不同的有限根。

\begin{proposition}\label{g4:prop:0204:root-cover}
	设 \(P(z,w)\in\mathbb C(z)[w]\) 首一、可分，次数为 \(n\ge1\)，并令有限集合 \(S\) 与 \(U\) 满足上述条件。则
	\[
		E_P=\{(z,w)\in U\times\mathbb C:P(z,w)=0\}
		\longrightarrow U,\qquad (z,w)\longmapsto z,
	\]
	是 \(n\) 叶全纯无分歧覆盖。按照穿孔圆盘幂模型补点，得到唯一的紧化 \(X_P\)，它是有限个紧黎曼曲面的不交并；投影延拓为到 \(\mathbb P^1\) 的全纯有限映射。坐标函数 \(z,w\) 均延拓为 \(X_P\) 上的亚纯函数。
\end{proposition}
\begin{proof}
	在 \(z_0\in U\) 上，全部根 \(w_1,\ldots,w_n\) 都是简单根。围绕它们各取一个小圆周 \(C_j\)，使相应闭圆盘两两不交且各只含这一根。由于 \(P\) 的系数在 \(z_0\) 附近全纯，缩小 \(z_0\) 的邻域后，\cref{g4:thm:0102:rouche}使每个 \(C_j\) 内始终恰有一个根。由\cref{g4:thm:0102:residue}，这个根可以写成
	\[
		w_j(z)=\frac1{2\pi\mathrm{i}}\int_{C_j}
		\frac{\xi\,P_w(z,\xi)}{P(z,\xi)}\,\mathrm{d}\xi.
	\]
	积分圆周上的分母不为零，并在缩小的闭参数圆盘上一致远离零，故可对 \(z\) 在积分号下求导；于是 \(w_j(z)\) 全纯。这些根彼此不同，且已给出次数 \(n\) 方程的全部根，所以得到覆盖的局部形式。

	在每个 \(b\in S\) 上方，应用\cref{g4:lem:0202:punctured-disk-cover}给各连通分支补点；与\cref{g4:thm:0204:monodromy-construction}证明中的紧性论证相同，有限叶性保证补成的空间紧。每个连通分支至少遇到一次 \(n\) 点纤维，所以分支数有限。补点的坐标与等价延拓唯一，故这一紧化在保持原覆盖的双全纯等价下唯一。

	剩下的实质问题是 \(w\) 在新点处会不会出现本性奇点。对任意首一复系数多项式
	\[
		w^n+c_1w^{n-1}+\cdots+c_n=0,
	\]
	每个根都满足
	\begin{equation}\label{g4:eq:0204:root-bound}
		|w|\le1+\max_{1\le j\le n}|c_j|.
	\end{equation}
	事实上，记右边的最大值部分为 \(A\)。若 \(|w|>1+A\)，则
	\[
		\left|\sum_{j=1}^n c_jw^{n-j}\right|
		\le A\sum_{j=1}^n|w|^{n-j}
		=A\frac{|w|^n-1}{|w|-1}<|w|^n,
	\]
	与方程矛盾。

	在 \(b\) 附近取目标坐标 \(u\)，使 \(u(b)=0\)；若 \(b=\infty\)，取 \(u=1/z\)。有理系数 \(a_j\) 在 \(u=0\) 至多有极点，因此存在整数 \(N\ge0\) 及常数 \(C\)，使穿孔小圆盘中
	\[
		|a_j(z(u))|\le C|u|^{-N}\qquad(1\le j\le n).
	\]
	由\eqref{g4:eq:0204:root-bound}，每个根满足 \(|w|\le C'|u|^{-N}\)。在某个新点的坐标 \(t\) 中，投影为 \(u=t^e\)，故
	\[
		|w(t)|\le C'|t|^{-eN}.
	\]
	函数 \(t^{eN}w(t)\) 有界且在穿孔圆盘上全纯，由\cref{g4:prop:0102:singularity-tests}延拓到 \(0\)。所以 \(w\) 在新点至多有极点。函数 \(z\) 的延拓更直接：在有限 \(b\) 上有 \(z=b+t^e\)，在无穷远上有 \(z=t^{-e}\)。因此二者均亚纯。
\end{proof}

这个根的增长界还把单值的轨道变成有理系数因子。它是“覆盖连通”与“多项式不可约”之间的关键一步。

\begin{theorem}\label{g4:thm:0204:irreducible-connected}
	设 \(P(z,w)\in\mathbb C(z)[w]\) 是次数 \(n\ge1\) 的首一可分多项式。取包含无穷远和全部系数极点的有限集合 \(S\)，使 \(z\notin S\) 时 \(P(z,w)\) 有 \(n\) 个不同根，令 \(E_P\to U=\mathbb P^1\setminus S\) 为根覆盖，\(X_P\) 为上述补点所得的紧化。以下条件等价：
	\begin{equivalences}
		\item \(P\) 在 \(\mathbb C(z)[w]\) 中不可约；
		\item 根覆盖 \(E_P\) 连通；
		\item 单值作用在任一基点的 \(n\) 个根上传递；
		\item \(X_P\) 连通。
	\end{equivalences}
	更精确地，\(P\) 的各个首一不可约因子，与 \(E_P\) 的连通分支、单值轨道及 \(X_P\) 的连通分支分别一一对应。
\end{theorem}
\begin{proof}
	每个覆盖分支都遇到基点纤维；同一纤维中的两点能被道路连接，恰好表示某条底空间环路的提升把前一点送到后一点。因此连通分支与单值轨道对应。每个新补点都有一个圆盘邻域，其穿孔部分连通，只属于原覆盖的一个分支，所以补点不会把两个分支合在一起。于是后三个条件等价。

	设单值不传递，在基点纤维中选择一个大小为 \(k\) 的轨道 \(O\)，其中 \(1\le k<n\)。在基点附近把根写成全纯函数 \(w_1(z),\ldots,w_n(z)\)，并组成
	\[
		Q_O(z,w)=\prod_{j\in O}(w-w_j(z)).
	\]
	其系数是这一组根的初等对称函数。沿任意路径继续根，可以继续这些系数；沿闭路返回时，轨道 \(O\) 被整体保持，根只在其中置换，所以每个系数返回原函数芽。由此它们在 \(U\) 上粘成单值全纯函数。

	还须证明这些系数是有理函数。在任意 \(b\in S\) 的穿孔坐标邻域中，上一命题的根界给出 \(|w_j|\le C|u|^{-N}\)。\(Q_O\) 的第 \(\ell\) 个系数是 \(\binom{k}{\ell}\) 个 \(\ell\) 重乘积的和，因此
	\[
		|q_\ell(u)|\le\binom{k}{\ell}C^\ell|u|^{-N\ell}.
	\]
	所以 \(u^{N\ell}q_\ell(u)\) 在穿孔圆盘上有界；由\cref{g4:prop:0102:singularity-tests}，它全纯延拓到原点。这说明每个 \(q_\ell\) 在包括无穷远在内的全部例外点至多有极点，即成为球面上的亚纯函数。

	球面上的亚纯函数必为有理函数，理由如下。它的极点局部有限，而球面紧，所以只有有限多个。减去所有有限极点的有限主部后，得到一个整函数 \(h(z)\)。由于无穷远至多为极点，存在整数 \(M\ge0\)，使 \(|h(z)|\le C(1+|z|)^M\)。对任意固定 \(a\)，在半径 \(R\) 的圆周上使用 Cauchy 估计，有
	\[
		|h^{(M+1)}(a)|
		\le (M+1)!\,\frac{C(1+|a|+R)^M}{R^{M+1}}
		\longrightarrow0.
	\]
	故 \(h\) 是次数不超过 \(M\) 的多项式。加回各有限主部，就得到有理函数。应用于 \(q_\ell\)，便知 \(Q_O\in\mathbb C(z)[w]\)。对其余轨道的根作同样构造，得到正次数首一多项式 \(R\)，并且 \(P=Q_OR\)。因此不传递必导致可约。

	反过来，设 \(P=QR\)，其中 \(Q,R\in\mathbb C(z)[w]\) 首一且均有正次数。先核对两因子的系数在 \(U\) 上没有极点。在一个 \(z_0\in U\) 的小闭圆盘内，\(P\) 的系数有界，所以全部根由\eqref{g4:eq:0204:root-bound}统一有界。离开 \(Q,R\) 系数的有限个可能极点时，每个因子的根都是 \(P\) 的一部分根，它们的对称函数因而有同一个局部上界；这些系数的可能极点遂由可去奇点判据消去。于是 \(Q,R\) 的系数都在 \(U\) 全纯。每个纤维中二者都至少有一个根，并且没有公共根，否则 \(P\) 有重根。故
	\[
		E_P=\{Q=0\}\ \sqcup\ \{R=0\}
	\]
	是两个非空闭子集的不交并；每个子集的补集也闭，所以它们都在 \(E_P\) 中开。根覆盖因此不连通。

	最后，对每个轨道作出的 \(Q_O\) 必不可约；若它可约，刚才的论证又把这个轨道分成两个非空不变部分，与轨道的定义矛盾。各轨道的积给出 \(P\)，所以得到所述逐因子对应。
\end{proof}

这里的连通对象是由不同根继续并补成的光滑曲面。原始平面解集可能在奇点处把不同分支碰到一起。例如
\[
	w^2-z^2=(w-z)(w+z)
\]
的仿射零点集是两条相交直线，作为集合连通；但在 \(z\ne0\) 上，根 \(w=z\) 与 \(w=-z\) 各自单值，根覆盖有两个分支。补点时每个分支各补自己的点，得到两张不交的射影直线。相反，\(w^2-z^3\) 在 \(\mathbb C(z)[w]\) 中不可约：若可约，就有有理函数 \(r(z)\) 满足 \(r(z)^2=z^3\)，而左边在 \(0\) 的阶数为偶数，右边为 \(3\)，矛盾。它的光滑参数为
\[
	z=t^2,\qquad w=t^3.
\]
在 \(z\ne0\) 处逆参数为 \(t=w/z\)，所以参数化确实恢复根覆盖；\(t=0\) 补一个点，\(t=\infty\) 再补一个点，得到一张连通曲面。方程的奇点、不可约性和覆盖的分歧分别在这里承担不同的信息。

\ProseParagraphBreak

分支位置与各点的分歧指数给出局部单值的循环类型，但没有给出各个置换如何共同作用于同一纤维。下面的例子把这一区别落实为可核验的数据。

\begin{example}\label{g4:exm:0204:same-branch-profile}
	固定球面上四个不同的点 \(b_1,b_2,b_3,b_4\)、基点及满足乘积关系的标准环路。取三点纤维，考虑两个元组
	\[
		\begin{aligned}
			\boldsymbol\sigma&=\bigl((12),(12),(23),(23)\bigr),\\
			\boldsymbol\tau&=\bigl((12),(23),(23),(12)\bigr).
		\end{aligned}
	\]
	两个元组按规定次序的复合均为恒等；它们都生成由 \((12),(23)\) 生成的传递群。因此\cref{g4:thm:0204:monodromy-construction}分别给出连通三次覆盖。每个 \(b_j\) 上方，两种覆盖都各有一个指数为 \(2\) 的点和一个指数为 \(1\) 的点；分支值恰好是同一组四点，所有局部分歧数据完全相同。

	这两个元组却不同时共轭。在 \(\boldsymbol\sigma\) 中，第一项等于第二项；在 \(\boldsymbol\tau\) 中，第一项不等于第二项，而同时共轭保持这种相等关系。因此两覆盖在固定目标投影下不等价。它们的源曲面亏格却相同，因为 Hurwitz 公式都给出
	\[
		2g-2=3(-2)+4=-2,\qquad g=0.
	\]
\end{example}

完整的单值元组因而同时记录局部循环和这些循环在同一组片上的配合；传递性检测连通性，循环长度给出补点与分歧指数，Hurwitz 公式再求出亏格。对有理系数代数方程，单值轨道还能借根增长界转成有理系数因子，把这一几何分解准确地接回不可约分解。
% 纲要条目排为编号小节标题，并进入目录与书签；说明段仍保留于纲要注释中。
